\subsection*{The physical endpoint of the complete Weil eigenfunctions}

The endpoint of an actual eigenfunction is not determined by ordinary
norm convergence of finite eigenvectors.  For the complete operator it
can instead be determined from the singular archimedean kernel.
Throughout this subsection $\Omega=(0,L)$, physical functions are
extended by zero outside $\Omega$, and Fourier frequency is normalized
as in $U_n(x)=L^{-1/2}e^{2\pi inx/L}$.

\begin{proposition}[Restricted logarithmic Laplacian decomposition]
\label{prop:v118-log-dirichlet}
Let $\mathcal A=\tfrac12L_\Delta^{\rm Dir}$ be the restricted,
exterior-Dirichlet logarithmic Laplacian on $\Omega$, whose full-line
symbol is $\log|t|$.  This is not the spectral logarithm of the
Dirichlet Laplacian.  The actual closed Weil operator satisfies
\begin{equation}\label{eq:v118-log-decomposition}
 \mathsf W_\lambda=\mathcal A+\mathcal K_\lambda,
 \qquad \mathcal D(\mathsf W_\lambda)=\mathcal D(\mathcal A),
\end{equation}
where, writing $k(y)=\rho(y)-(2y)^{-1}$, $a_m=\log m$, and
$w_m=\Lambda(m)/\sqrt m$,
\begin{align}
 (\mathcal K_\lambda f)(x)={}&-\log(2\pi)f(x)
       -\int_0^L k(|x-y|)f(y)\,dy\notag\\
 &-\sum_{1<m<\lambda^2}w_m
       \{\mathbf1_{x+a_m<L}f(x+a_m)
                   +\mathbf1_{x>a_m}f(x-a_m)\}\notag\\
 &+2\int_0^L\cosh((x-y)/2)f(y)\,dy.
 \label{eq:v118-bounded-physical-part}
\end{align}
The operator $\mathcal K_\lambda$ is real and self-adjoint on $L^2$.
For every $1\le p\le\infty$ it is bounded on $L^p(\Omega)$, with
\begin{equation}\label{eq:v118-K-bound}
 \|\mathcal K_\lambda\|_{p\to p}\le C_\lambda
 :=\log(2\pi)+2\int_0^L|k(y)|\,dy
                  +2P_\lambda+4\sinh(L/2).
\end{equation}
\end{proposition}
\begin{proof}
In dimension one the integral normalization of the logarithmic
Laplacian~\cite[Theorem~1.1]{ChenWethLog} gives
\[
 \mathcal A f(x)=\int_0^\infty
 \frac{2f(x)\mathbf1_{y<1}-f(x+y)-f(x-y)}{2y}\,dy
                   -\gamma_{\rm E}f(x)
\]
on compactly supported smooth tests.  The actual negative-archimedean
part of the Weil operator is
\[
 \int_0^\infty\rho(y)
 \{2e^{-y/2}f(x)-f(x+y)-f(x-y)\}\,dy
             -\{\log(4\pi)+\gamma_{\rm E}\}f(x).
\]
Its Fourier diagonal is precisely the negative of
\eqref{eq:v114-arch-diagonal}; the diagonal integral over $y>L$
equals $-\log\tanh(L/2)$.  Since
\[
 \int_0^\infty\left\{\frac1{\sinh y}
                -\frac{\mathbf1_{y<1}}y\right\}dy=\log2,
\]
subtracting $\mathcal A$ leaves $-\log(2\pi)I$ and the kernel $-k$.
The integral follows by evaluating $\log\tanh(y/2)$ with a positive
lower cutoff and then taking its limit.  In particular the constant
part of the logarithmic diagonal has been retained.
The prime shifts and the pole kernel $2\cosh((x-y)/2)$ give the
remaining terms of \eqref{eq:v118-bounded-physical-part}.
The latter kernel is the difference of the centered rank-one
$2\cosh\otimes\cosh$ and $2\sinh\otimes\sinh$ terms, so no pole
sign has been removed.

The function $k$ extends continuously to zero, with $k(0)=1/4$.
The integral-kernel row and column bounds, and the contraction of
each truncated translation, prove \eqref{eq:v118-K-bound}.
For the pole kernel the largest row integral is
$4\sinh(L/2)$.  Symmetry proves self-adjointness.

It remains to identify the closed forms, without imposing an
unproved endpoint condition.  Smooth compactly supported functions
are a form core for the restricted logarithmic
Laplacian~\cite[Theorem~3.1]{ChenWethLog}.  They are also a form core
for the operator in Proposition~\ref{prop:v114-weil-core}.
Indeed, approximate each finite periodic Fourier sum $f$ by smooth
cutoffs that remove endpoint strips of width $\epsilon$.
The discarded periodic function has $L^2$ norm $O(\epsilon^{1/2})$
and $H^1$ norm $O(\epsilon^{-1/2})$, hence $H^s$ norm
$O(\epsilon^{1/2-s})$ for any fixed $0<s<1/2$.
This tends to zero in the logarithmically weighted form norm.
The two closed lower-bounded forms therefore agree on a common form
core.  Bounded perturbation by $\mathcal K_\lambda$ proves
\eqref{eq:v118-log-decomposition}, including its operator-domain assertion.
\end{proof}

\begin{theorem}[Zero boundary trace of every complete Weil eigenfunction]
\label{thm:v118-eigen-endpoint-zero}
Fix $\lambda>1$.  Every $L^2$ eigenfunction
$\mathsf W_\lambda u=\mu u$ has a bounded representative whose zero
extension is continuous on $\mathbb R$.  With
\[
 \ell(r)=\frac1{|\log(\min\{r,0.1\})|},\qquad
 d(x)=\min\{x,L-x\},
\]
there is a finite constant $C$ such that
\begin{equation}\label{eq:v118-eigen-boundary-decay}
 |u(x)|\le C\sqrt{\ell(d(x))}\quad(0<x<L).
\end{equation}
In particular its actual physical endpoints satisfy $u(0)=u(L)=0$.
No uniform bound for $C$ as $\lambda$ grows is asserted.
\end{theorem}
\begin{proof}
First we prove boundedness, without assuming a positive Sobolev gain.
For $0<\epsilon<\min\{1,L\}$ let $\mathcal A_\epsilon$ be the
nonnegative killed small-jump generator
\[
 \mathcal A_\epsilon f(x)=\int_0^\epsilon
       \frac{2f(x)-f(x+y)-f(x-y)}{2y}\,dy
\]
in its closed Dirichlet-form realization, and put
\[
 J_\epsilon f(x)=\int_{\substack{y\in\Omega\\|x-y|\ge\epsilon}}
                    \frac{f(y)}{2|x-y|}\,dy.
\]
There is an exact decomposition
\[
 \mathcal A=\mathcal A_\epsilon+
       \{\log(1/\epsilon)-\gamma_{\rm E}\}I-J_\epsilon.
\]
The small-jump form is Markovian, including its nonnegative exterior
killing term.  Consequently its resolvent is consistent on $L^2$
and $L^\infty$, and
$\|(\mathcal A_\epsilon+aI)^{-1}\|_{p\to p}\le a^{-1}$ for
$a>0$, $p=2,\infty$.  One may obtain this directly by first cutting
off jumps smaller than a positive number, using the contraction
resolvent of that finite-activity killed jump generator, and then
passing to the increasing closed forms.  Also Cauchy--Schwarz gives
\[
 \|J_\epsilon\|_{2\to\infty}\le(2\epsilon)^{-1/2}.
\]
Choose $\epsilon$ so that
$a=\log(1/\epsilon)-\gamma_{\rm E}-\mu>C_\lambda$.
The eigen equation becomes
\[
 (\mathcal A_\epsilon+aI+\mathcal K_\lambda)u=J_\epsilon u.
\]
The inverse on the left exists on both $L^2$ and $L^\infty$ by the
Neumann series for
$[I+(\mathcal A_\epsilon+aI)^{-1}\mathcal K_\lambda]^{-1}
 (\mathcal A_\epsilon+aI)^{-1}$, with norm at most
$(a-C_\lambda)^{-1}$.  Since $J_\epsilon u\in L^2\cap L^\infty$,
the two series agree.  Uniqueness in $L^2$ identifies the bounded
solution with $u$, and proves
\[
 \|u\|_\infty\le
   \frac{\|u\|_2}{(a-C_\lambda)\sqrt{2\epsilon}}.
\]
Thus boundedness has not been inferred from the logarithmic operator
domain alone.

By Proposition~\ref{prop:v118-log-dirichlet}, the zero extension lies
in the logarithmic Dirichlet form space and satisfies
\[
 L_\Delta u=2(\mu u-\mathcal K_\lambda u)\in L^\infty(\Omega),
 \qquad u=0\quad\hbox{outside }\Omega.
\]
The interval satisfies an exterior uniform sphere condition.
The boundary regularity theorem of Hern\'andez-Santamar\'ia,
L\'opez R\'ios and Salda\~na~\cite[Theorem~1.1]{HLSLogBoundary}
therefore gives continuity of the zero extension and
\eqref{eq:v118-eigen-boundary-decay}.  For complex $u$, apply the
real theorem to its real and imaginary parts.  This proves the endpoint
assertion for every eigenfunction, independently of its spectral order.
\end{proof}

\begin{corollary}[Filtered endpoint decay and the nonzero-source obstruction]
\label{cor:v118-ground-fejer}
For an eigenfunction as above and $N\ge4$, define its Fej\'er endpoint
in the actual normalized Fourier coordinates by
\[
 \sigma_{N-1}u(0)=L^{-1/2}
       \sum_{|n|<N}(1-|n|/N)\widehat u(n).
\]
Then
\begin{equation}\label{eq:v118-fejer-endpoint}
 |\sigma_{N-1}u(0)|
 \le C\sqrt{\ell(L/\sqrt N)}+\frac{\|u\|_\infty}{2\sqrt N}
 =O_{\lambda,\mu,u}((\log N)^{-1/2}).
\end{equation}
For the exact repaired source with $B_\lambda=p_\lambda(0)\ne0$,
any complete-operator eigenfunction $u$ and scalar $c$ satisfy
\begin{equation}\label{eq:v118-ground-source-mismatch}
 \frac{|p_\lambda(0)-cu(0)|}{|B_\lambda|}=1.
\end{equation}
\end{corollary}
\begin{proof}
The positive Fej\'er kernel in physical coordinates is
$L^{-1}F_N(2\pi x/L)$, where
$F_N(\theta)=N^{-1}\{\sin(N\theta/2)/\sin(\theta/2)\}^2$;
it has integral one.  Put $h=L/\sqrt N\le L/2$.
On $d(x)<h$, \eqref{eq:v118-eigen-boundary-decay} bounds the
convolution by $C\sqrt{\ell(h)}$.
On $d(x)\ge h$, $\sin(\pi x/L)\ge2d(x)/L$ gives
\[
 L^{-1}F_N(2\pi x/L)\le\frac{L}{4N d(x)^2},\qquad
 \int_{d(x)\ge h}L^{-1}F_N(2\pi x/L)\,dx
       \le\frac{L}{2Nh}=\frac1{2\sqrt N}.
\]
This proves \eqref{eq:v118-fejer-endpoint}.
The mismatch identity follows from $u(0)=0$.
\end{proof}

Theorem~\ref{thm:v118-eigen-endpoint-zero} applies in particular to
the simple even ground at $\lambda=3$.  It disproves a nonzero physical
endpoint comparison between the repaired source and the complete
operator's ground, while leaving their ordinary-norm overlap intact.
It does not identify the endpoint of a finite-compression eigenvector
on a joint growing-window and growing-cutoff diagonal.
Nor does continuity alone prove absolute Fourier convergence or
convergence of sharp Fourier endpoint sums: the quantitative statement
\eqref{eq:v118-fejer-endpoint} concerns the stated positive filter.
The finite endpoint and its graph-defect terms therefore remain
separate from the complete eigenfunction's trace.
